% Transcription — fonds Alexandre Grothendieck, shelfmark 43, batch 7 (pages 121-140).
% Established from the facsimile archives/batches/43/batch-07.pdf (PDF page k =
% archivists' page 120+k, no cover sheet), cut from a copy of 43.pdf fetched
% through the project's own relay
% (https://grothendieck-relay.onrender.com/source/43.pdf); tiles in
% archives/tiles/43/p121-p140. Per the skill transcribe-grothendieck. The folder
% is 207 pages in eleven batches, transcribed in parallel. Batch 5 was read for
% the folder's conventions, and batch 6 (committed during this pass, ending on
% the EGA III leaf of p. 120) was read and followed: \underline{\mathrm{Ext}}
% for his underlined Ext (doubly underlined where he doubles the stroke),
% \mathbb{G}_m, \gamma for the coefficient k-group, S\gamma = (D\Delta\gamma)(1),
% \mu_\infty, \check{\mu}_n for his mu with a caron, \hat{F} for the dual.
%
% Pass: Opus 5.5 (claude-opus-5-5), 2026-09-26 — first pass, unchecked against the pages by a human.
%
% Dating: the folder has its own inventory date, carried verbatim in \dating{}
% with no group sentence (the group « Jacobiennes » (43-44) is added by the
% skill only for undated folders). No page of this batch carries a date.
%
% Copies. No page of this batch can be identified as a photocopy. The notes
% are in ink (black, blue, blue-black) and pencil on original paper: pp. 122,
% 124, 135, 137, 139 are written on the blank backs of typed leaves whose
% typing shows through in mirror; p. 128 carries ink blots and blue ink over
% pencil strokes; p. 127 an ink blot. All are transcribed as originals.
%
% Pages skipped: 123, 125, 138, 140 — typed leaves (upside down for 123 and
% 125) of a French typescript on sheaves, headed « n° 239 », typed pages 84,
% 78, 88, 87 (Théorème 3 on soft sheaves of rings and partitions of unity;
% § 3 « Faisceaux mous », Définition 1, axioms (FM_I), (FM_II); Proposition 8
% and Lemme 2 on fine sheaves), with nothing in his hand; they are the other
% sides of pp. 122, 124, 137, 139. Page 126: a blank sheet.
%
% Typescripts. Page 136 is a typed leaf paginated « - 37 - », « n° 338 »
% (faithfully flat ring extensions, conditions (iv)-(v), the proof
% (i) => (iii) => (iv) => (ii) => (i), Proposition 9 begun), carrying one
% pencil marginal note of his: one French \note, the typed text not re-set,
% following batch 6's treatment of the EGA III leaves.
%
% Cover: page 121 is a tan folder sheet inscribed at the top right in his
% hand « Corps de classes et dualité des Courbes Algébriques » (the last two
% words underlined): no \page marker, his words become the section heading,
% with a \note. It opens the run pp. 127-139 on curves.
%
% Pages summarised: none. Pages 130 and 133 are short cancelled headings
% (a « V », then « Schémas de Hilbert, de Picard » / « Fidèle platitude,
% Descente fid. plate, Applications »), given in full. Page 128 (cancelled by
% long strokes; commutative algebra, dimension of localised polynomial rings)
% and the prose of pp. 127, 129, 131, 134 are fast drafting and carry a heavy
% apparatus; the scattered formula sheets 122, 124, 135, 137, 139 are given
% block by block, their layout said in \note{}s; the table of p. 132 is set
% as a list of its cells.
%
% His pagination: none. Typed page numbers on the skipped leaves: « n° 239 »
% pp. 84 (p. 123), 78 (p. 125), 88 (p. 138), 87 (p. 140); « n° 338 » p. 37
% (p. 136). A Roman « V » heads pp. 130 and 133 (on 130 followed by a struck
% numeral). Numbered runs: a), b) (*), c) on p. 131 (the pairing to be
% defined); Th. (i), (ii) on p. 129; (i), (ii), Cor. (three times) on p. 128;
% the systems of three isomorphisms in the box of p. 135. No internal page
% reference.
%
% What the batch is about. Duality for curves over a field, as the global
% counterpart of the local duality of batches 5-6. P. 122: local
% cohomology H^i_Y(X, M), H^i(X, M), H^i(U, M) against Ext^i(X; M, A),
% Ext^i_Y(X; M, A), Ext^i(U; M, A), « Ext_m(M, A) =? D(M) », the pairing of
% Ext^i_Y(X; M, N) with Ext^{n-i}(X; N, M) along the two local-cohomology
% sequences, and the formal completions X_{/Y} -> X_{/Z} for {a} ⊂ Y ⊂ Z.
% P. 124: Ext^i_Y(X, F, G) = lim Ext^i(F/J^n F, G), dual to lim Ext^{r-i}_a,
% the spectral sequences H^p_a(X, Ext^q), the passage to the completion,
% then a formal scheme 𝔛, T ⊂ Z ⊂ 𝔛 and H^i_T(𝔛, F) -> H^i_T(𝔛, F_{/Z}).
% P. 127: C a complete nonsingular curve over k, C-groups (Weil sheaves),
% f_* : Gr_C -> Gr_k -> QuGr_k, its derived functors H^i_f and Ext^i_f and
% the Leray-type spectral sequences. P. 128 (cancelled): dim B_m = dim A + n -
% deg tr k(m)/k(m) for B = A[t_1..t_n], and a corollary on X = Y[t_1..t_n].
% P. 129: the construction of the duality map E_f(G, S gamma) -> E_f(Rf_*G,
% gamma) via a trace Rf_* S gamma -> gamma, and a Theorem (i)-(ii) (G of
% finite type, then without unipotent (?) component). Pp. 130, 133:
% cancelled chapter headings. P. 131: the Kummer sequence, H^i_f(G_m) = 0 for
% i >= 2, 0 -> H^1(k, G_m) -> H^1(C, G_m) -> Pic(C/k) -> Br(k), S gamma =
% (D Delta gamma)(1), and the maps a), b), c). P. 132: a table of the duality
% for G = mu_n, Z/nZ, mu_p, Z/pZ, alpha_p, G_m, Z (with the Jacobian J,
% Frobenius-fixed and co-fixed parts of H^1(C, O_C), self-duality of the
% Jacobian), then H*(C, gamma) ≃ Ext*(Rf_*(G_m)(1), gamma) and H^i(C, Q/Z).
% P. 134: reduction to a sheaf concentrated at a point x, T(G) ≃ G_x, and the
% boxed H^i(U; F^) ≃ Ext^{i-2}(Rf_*(i_!F_U), Q/Z). P. 135: F = i_*(G) on the
% spectrum S of a complete DVR, H^2_x(F) = H^1(K, F), the triangle
% H_x(F) -> H_X(F) -> H_U(F) and three duality isomorphisms with Q/Z.
% P. 137: the two long exact sequences of H_{Y/k}, H_{S/k}, H_{U/k} for G
% and G^ set face to face. P. 139: Ext_{U/k}(G, gamma~) ≃ Ext_k(N(G), gamma)
% (N the Néron functor), 0 -> G_U -> G -> G^Y -> 0, with « ??? ».
%
% Notation fixed once for the batch: C the curve, f : C -> Spec k; \mathcal{N}
% for his curly N (the Néron functor of batch 5), written N^{\cdot} where he
% puts a dot; \mathfrak{X} for his gothic X (p. 124); \tilde{\gamma} as he
% writes it (p. 139); \mathrm{Pic}, \mathrm{Br}; \underline{\underline{\mathrm{Ext}}}
% for his doubly underlined Ext (p. 132).
%
% Readings to check first: the whole prose of pp. 127, 129 and 134; the
% cancelled p. 128; the margin of p. 132 and the three glosses of its right
% column; the twist « (1) » or « (2) » in S gamma on p. 131; the exponent of
% the right-hand Ext in a) on p. 131; the last term of p. 139.
% Apparatus: 217 \ill{} and 331 \uncertain{} in the body.

\documentclass[11pt,a4paper]{article}
% Le préambule commun à toutes les transcriptions du fonds Grothendieck.
%
% Il tient en une idée : **l'incertitude se compose, elle ne se résout pas.**
% Une transcription qui lisse un mot douteux a détruit la seule chose pour
% laquelle on la faisait — un lecteur doit pouvoir savoir, à chaque endroit,
% ce qui était sur le feuillet et ce qui a été ajouté. Les six macros qui
% suivent sont donc l'appareil critique, et non de la décoration.
%
% Les mêmes commandes sont reconnues par `scripts/render.mjs`, qui produit la
% vue de lecture du site. Une macro ajoutée ici doit l'être aussi là-bas,
% faute de quoi le rendu échouera bruyamment — ce qui est voulu : un rendu
% silencieusement faux est pire qu'un rendu absent.

\ProvidesPackage{grothendieck}[2026/08/08 Transcriptions du fonds Grothendieck]

\usepackage{amsmath,amssymb,amsthm}
\usepackage{mathrsfs}
\usepackage{stmaryrd}
% Les diagrammes commutatifs sont partout dans ce fonds : c'est la langue
% même de la géométrie algébrique de Grothendieck, et une transcription qui
% les rendrait en prose ne transcrirait rien.
\usepackage{tikz-cd}
% Français par défaut — c'est la langue des feuillets — mais l'anglais est
% chargé aussi : la traduction et le résumé sont des documents anglais, et la
% typographie française (espaces avant les deux-points) y serait une faute.
% Ces documents basculent par \selectlanguage{english} en tête de corps.
\usepackage[english,french]{babel}
\usepackage{fontspec}
\usepackage{geometry}
\geometry{a4paper,margin=25mm,marginparwidth=28mm}
\usepackage{marginnote}
\usepackage{xcolor}
% ulem, not soul. `soul`'s \st and \ul are built on a character-by-character
% scan that XeLaTeX's font machinery breaks: the document fails with an
% undefined control sequence rather than degrading. ulem does the same two jobs
% and is Unicode-engine safe. `normalem` stops it hijacking \emph, which the
% transcriptions use for Grothendieck's own emphasis.
\usepackage[normalem]{ulem}
\usepackage{hyperref}
\hypersetup{colorlinks=true,linkcolor=black,urlcolor=black,pdfborder={0 0 0}}

% --- Métadonnées du lot -----------------------------------------------------
% Reprises telles quelles par le script de rendu, qui les affiche en tête de
% la vue de lecture. Elles fixent ce que le fichier prétend couvrir : sans
% elles, une transcription flotte sans point d'attache dans un fonds de seize
% mille feuillets.

\newcommand{\folder}[1]{\gdef\grothFolder{#1}}
\newcommand{\batch}[1]{\gdef\grothBatch{#1}}
\newcommand{\leaves}[2]{\gdef\grothFirst{#1}\gdef\grothLast{#2}}
\newcommand{\foldertitle}[1]{\gdef\grothFoldertitle{#1}}
\gdef\grothFolder{?}\gdef\grothBatch{?}\gdef\grothFirst{?}\gdef\grothLast{?}\gdef\grothFoldertitle{}

% --- Le résumé --------------------------------------------------------------
%
% Il ouvre la lecture modernisée, sous le titre « Résumé », et il est composé
% en petit. Ce n'est pas un ornement : c'est le seul passage du document qui
% ne soit pas des mathématiques, et un lecteur doit pouvoir le reconnaître
% avant de l'avoir lu — puis le sauter s'il connaît déjà le sujet, ou s'y
% arrêter s'il ne le connaît pas. Le filet qui le suit marque la reprise du
% corps.

\newenvironment{resume}
  {\small\setlength{\parskip}{0.45em}}
  {\par\medskip\hrule\medskip}

% --- Mots-clés ---------------------------------------------------------------
%
% En anglais, en clôture du résumé de la lecture modernisée : le vocabulaire
% moderne sous lequel chercher ce que les feuillets construisent. Le manifeste
% du site (`npm run manifest`) les extrait pour étiqueter la cote entière —
% cette ligne est la source unique des tags, il n'y a pas de fichier à côté.
\newcommand{\keywords}[1]{\par\smallskip\noindent
  {\footnotesize\textsc{Keywords} --- \textit{#1}}\par}

% --- L'appareil critique ----------------------------------------------------

% Le numéro de feuillet, en marge. C'est l'ancre qui permet de régler un
% désaccord sur une formule en regardant *un* feuillet plutôt qu'une chemise
% de six cents. Le site s'en sert aussi pour tourner les pages du fac-similé
% quand on fait défiler la transcription.
\newcommand{\leaf}[1]{%
  \marginnote{\footnotesize\color{gray!70}\textsf{#1}}%
  \label{leaf:#1}\ignorespaces}

% Illisible : on ne devine pas. Un mot inventé qui se lit comme les autres est
% le pire résultat possible.
\newcommand{\ill}{\textcolor{red!60!black}{[\,\ldots\,]}}

% Lecture douteuse : le mot est donné, et signalé comme incertain.
\newcommand{\uncertain}[1]{\uline{#1}}

% Ajout de l'éditeur — une lettre manquante, un mot sous-entendu.
\newcommand{\add}[1]{\textcolor{blue!50!black}{[#1]}}

% Biffé par Grothendieck lui-même. Ce qu'il a rayé fait partie du document :
% une rature dit souvent où la pensée a changé de direction.
\newcommand{\struck}[1]{\sout{#1}}

% Note du transcripteur, sur l'état matériel du feuillet.
\newcommand{\note}[1]{\par\smallskip{\small\color{gray!60!black}\textit{[#1]}}\par\smallskip}

% Note marginale de Grothendieck, distincte du corps du texte.
\newcommand{\marginal}[1]{\par\smallskip\noindent\hspace{1em}%
  {\small\color{gray!60!black}#1}\par\smallskip}

% --- Titre ------------------------------------------------------------------

\AtBeginDocument{%
  \begin{center}
    {\large\bfseries Fonds Alexandre Grothendieck --- cote n\textsuperscript{o}~\grothFolder}\\[2pt]
    {\itshape\grothFoldertitle}\\[4pt]
    {\small feuillets \grothFirst--\grothLast\ (lot \grothBatch)}\\[2pt]
    {\footnotesize\color{gray!60!black}Transcription établie d'après le fac-similé numérisé
     par l'Université de Montpellier.\\
     Les lectures incertaines sont soulignées, les illisibles notées [\,\ldots\,],
     les ajouts entre crochets.}
  \end{center}
  \vspace{1em}}

\endinput


\folder{43}
\batch{7}
\pages{121}{140}
\dating{[à partir de 1961-vers 1968]}
\watermark{Édition de démonstration}
\foldertitle{Dualité des groupes. Jacobiennes généralisées etc : notes et copies de notes manuscrites (s.d.)}

\begin{document}

\section*{Corps de classes et dualité des Courbes Algébriques}
\note{titre inscrit de sa main en haut à droite d'une chemise de papier
kraft (page 121), « Courbes Algébriques » souligné, sans autre contenu ;
les notes sur les courbes suivent aux pages 127 et suivantes.}

\page{122}
\note{feuillet de formules éparses, écrit au dos d'une page dactylographiée
qui transparaît. En haut à gauche, un encadré à deux colonnes, mises ici
face à face.}
\[
\begin{array}{ll}
H^{i}_{Y}(X, M) & \mathrm{Ext}^{i}(X; M, A) \\
H^{i}(X, M) & \mathrm{Ext}^{i}_{Y}(X; M, A) \\
H^{i}(U, M) & \mathrm{Ext}^{i}(U; M, A)
\end{array}
\]
\note{dans la colonne de droite, les indices et le premier argument sont
surchargés : la première ligne semble porter un $Y$ biffé sous le $X$, la
seconde un $Y$ repassé en indice.}
\[
\mathrm{Ext}^{i}_{\mathfrak{m}}(M, A) \overset{?}{=} D(M)
\]
\[
0 \to M' \to M \to M'' \to 0
\]
\[
\ill{}\ (M, A)
\]
\note{au-dessous de la suite exacte, une ligne courte dont le début est
illisible, suivie de « $(M, A)$ ». Un long trait double oblique sépare cette
partie gauche du reste ; au milieu, un grand « $A$ » isolé.}

\[
\begin{array}{ccc}
\mathrm{Ext}^{i}_{Y}(X; M, N) & \cdot & \mathrm{Ext}^{n-i}(X; N, M) \\
\downarrow & & \uparrow \\
\mathrm{Ext}^{i}(X; M, N) & & \mathrm{Ext}^{n-i}_{Y}(X; N, M) \\
\downarrow & & \uparrow \\
\mathrm{Ext}^{i}(U; M, N) & & \mathrm{Ext}^{n-i}(U; N, M)
\end{array}
\]
\note{un point entre les deux colonnes de la première ligne marque
l'accouplement. L'exposant de droite est lu $n-i$ ; au premier terme de
gauche, le $X$ est surchargé.}

\[
\mathrm{Ext}^{i}_{Y}(X; F, G) \to \mathrm{Ext}^{i}_{Z}(X; F, G) \to \mathrm{Ext}^{i}_{Z-Y}(X; F, G)
\]
\note{à gauche de cette ligne, un petit cercle avec un point au centre.}
\[
\begin{split}
\mathrm{Ext}^{r-i}(X_{/Y} - Y, G_{/Y}, F_{/Y}) &\overset{?}{\longleftarrow} \mathrm{Ext}^{r-i}_{a}(X_{/Y}; \hat{G}_{/Y}, \hat{F}_{/Y}) \longleftarrow \mathrm{Ext}^{r-i}_{a}(X_{/Z}; G_{/Z}, F_{/Z}) \\
&\overset{?}{\longleftarrow} \mathrm{Ext}^{r-i-1}(X_{/Z} - Y, G_{/Z}, F_{/Z})
\end{split}
\]
\note{la ligne est écrite d'un seul tenant ; les indices « $/Y$ »,
« $/Z$ » sont des lectures de ses barres obliques. Au deuxième terme, les
chapeaux sont surchargés et un « $\approx$ ? » est écrit dessous ; au
dernier terme, l'exposant est lu $r-i-1$.}
\[
X_{/Y} \longrightarrow X_{/Z} \qquad Y \subset Z, \quad Z \supset Y \supset \lbrace a \rbrace
\]
\note{sous $X_{/Z}$, la chaîne d'inclusions est écrite verticalement :
$Z \supset Y \supset \lbrace a \rbrace$ ; « $Y \subset Z$ » est à gauche.}
\[
\ldots \to H^{\cdot}_{a}(X_{/Z}; G_{/Z}, F_{/Z})
\]
\note{le symbole après la flèche est peu lisible ; il porte l'indice $a$.}

\page{124}
\note{au dos d'une page dactylographiée qui transparaît.}
\[
\mathrm{Ext}^{i}_{Y}(X, F, G) = \varinjlim_{n} \mathrm{Ext}^{i}(F/J^{n}F, G)
\]
\[
\simeq \varprojlim_{n} \mathrm{Ext}^{r-i}_{a}(X; G, F/J^{n}F) \Longleftarrow H^{p}_{a}\bigl(X, \underline{\mathrm{Ext}}^{q}_{\mathcal{O}_{X}}(G, F_{n})\bigr)
\]
\[
\begin{array}{c}
\uparrow{\scriptstyle \wr} \\
\varprojlim_{n} \mathrm{Ext}^{r-i}_{a}(\hat{X}, \hat{G}, \hat{F}_{n}) \Longleftarrow H^{p}_{a}\bigl(\hat{X}, \underline{\mathrm{Ext}}^{q}_{\mathcal{O}_{\hat{X}}}(\hat{G}, \hat{F}_{n})\bigr) \\
\uparrow \\
\varprojlim \mathrm{Ext}^{r-i}_{a}(\hat{X}, \hat{G}, \hat{F}_{n})
\end{array}
\]
\note{« dual » est écrit dans la marge gauche en face du premier $\simeq$.
La dernière ligne reprend la précédente sans le $\varprojlim$ indicé ; la
flèche qui y monte est simple.}
\[
\mathrm{Ext}^{i}_{Y-Z}(X, F, G) = \varinjlim_{n} \mathrm{Ext}^{i}(X - Z, F/J^{n}F, G)
\]
\[
\simeq \varprojlim \ \ldots \qquad \varprojlim \mathrm{Ext}^{r-i}(\hat{X}, \hat{G}, \hat{F})
\]
\note{la seconde ligne est restée inachevée : un blanc entre le
$\varprojlim$ et le dernier terme.}
\[
H^{i}(X - Z, F
\]
$Z \subset Y$

$\mathfrak{X}$ \uncertain{schéma} \uncertain{formel}
\[
T \subset Z \subset \mathfrak{X}
\]
$Z$ \uncertain{partie} \uncertain{fermée} \ldots\ ;
$\mathfrak{X}_{/Z}$ \uncertain{complété} \uncertain{formel} \uncertain{le}
\uncertain{long} \uncertain{de} $Z$
\note{les deux gloses sont reliées à $\mathfrak{X}$ et à $\mathfrak{X}_{/Z}$
par des flèches.}
\[
H^{i}_{T}(\mathfrak{X}, F) \to H^{i}_{T}(\mathfrak{X}, F_{/Z}) \overset{?}{\longrightarrow} H^{i}(\mathfrak{X}
\]
\note{la ligne s'arrête au bord du texte.}

\page{127}
\struck{\uncertain{Schémas} \ill{}}

$C$ \uncertain{courbe} \uncertain{algébrique} [\uncertain{complète}
\uncertain{sans} \uncertain{singularités} \uncertain{sur} $k$]

On \uncertain{sait} \uncertain{définir} la \uncertain{notion} de
$C$-\uncertain{groupe} [\uncertain{faisceaux} de Weil \uncertain{pour} \ill{}
\uncertain{préschémas} de \struck{\ill{}} type \uncertain{fini} sur $C$,
\uncertain{par} la \uncertain{descente} \uncertain{fid.}\ \uncertain{plate}
\ldots]. \struck{\ill{}} Les $C$-groupes \uncertain{forment} une catégorie
\uncertain{abélienne} \uncertain{avec} \uncertain{assez}
d'\uncertain{injectifs}, d'où des $\mathrm{Ext}$ (globaux et locaux), des
$H^{i}(C, G) \simeq \mathrm{Ext}^{i}(C; \mathbb{Z}_{C}, G)$, etc.
\note{« $\simeq \mathrm{Ext}^{i}(C; \mathbb{Z}_{C}, G)$ » est écrit au-dessus
de la ligne, au-dessus de « etc. ».}

\uncertain{Sur} $k$, on a les $k$-\uncertain{groupes} [\uncertain{définis}
\uncertain{comme} les $C$-groupes ? \struck{\uncertain{Cela}} \ill{}
\uncertain{équivalent} aux \uncertain{faisceaux} de Weil
\add{\uncertain{définis}} \ill{} les $k$-\uncertain{ensembles}
\uncertain{algébriques} ?]
\note{le mot inséré au-dessus de la ligne, après « Weil », est peu lisible ;
le mot biffé devant « $k$-ensembles » aussi.}

\uncertain{On} \uncertain{a} un \uncertain{foncteur} $f_{*}$
\[
f_{*} \colon \mathrm{Gr}_{C} \longrightarrow \mathrm{Gr}_{k} \longrightarrow \mathrm{QuGr}_{k}
\]
[\uncertain{image} \uncertain{directe} par $f \colon C \to \mathrm{Spec}\, k$]
\uncertain{qui} \ill{} \uncertain{exact} \ill{} \struck{\ill{}} et
\uncertain{transforme} \uncertain{injectifs} en \uncertain{injectifs}, d'où
des \uncertain{foncteurs} \uncertain{dérivés}
\[
\underline{H}^{i}_{f}(G) = \struck{\ill{}}\ H^{i}(Rf_{*})(G)
\]
[\uncertain{où} $G$ \ill{} \uncertain{complexe} \ill{} \ill{} \ill{} \ldots\
\ill{} \uncertain{que} $f_{*}$ \ill{} \ill{} \ill{} \emph{\ill{}}] On
\uncertain{définit} de même
\[
\underline{\mathrm{Ext}}^{i}_{f}(F, G) = H^{i}\bigl(\underline{E}^{\cdot}_{f}(F, G)\bigr) = H^{i}\bigl(\bigl(R\,\underline{\mathrm{Hom}}_{f}(F, -)\bigr)(G)\bigr)
\]
D'\uncertain{où} \uncertain{comme} d'\uncertain{habitude} des
\uncertain{suites} \uncertain{spectrales}
\[
\mathrm{Ext}^{*}(C; F, G) \Longleftarrow H^{p}\bigl(k; \underline{\mathrm{Ext}}^{q}_{f}(F, G)\bigr)
\]
en \uncertain{particulier}
\[
H^{*}(C, G) \Longleftarrow H^{p}\bigl(k; \underline{H}^{q}_{f}(G)\bigr)
\]
\note{la suite spectrale des $\mathrm{Ext}$ est précédée, à gauche, d'un
$H^{*}_{f}$ \uncertain{biffé}.}

\page{128}
\note{page entière annulée par deux longs traits obliques ; lecture très
incertaine de la prose, qui est rapide.}
\marginal{\uncertain{passer} \ill{} \uncertain{N=} \ill{}}
\struck{\uncertain{Proposition}}
\note{la note de marge, en haut à gauche, est écrite en biais et reliée par
une flèche au début du (i).}

(i) \uncertain{Si} $A = k[t_{1}, \ldots, t_{n}]_{\mathfrak{m}}$, \ill{}
\ill{}, \ill{} \ill{} \ill{} \struck{$A =$ \ill{} $A$ \ill{} \ill{}}
\add{$A$ \uncertain{est} \uncertain{régulier}, \uncertain{de}
\uncertain{dimension}} $\dim A = n$ \struck{$A[t_{1}, \ldots, t_{n}]$}

\emph{Dém.} \uncertain{Par} \ill{} \ill{}, \ill{} \ill{} \ill{}
\uncertain{au} \uncertain{cas} \uncertain{où} $k(\mathfrak{m}) = k$, \ill{}
\ill{} \struck{\ill{}} \add{$A$-\ill{}} \ill{} \ill{} \ill{}
\uncertain{l'idéal} \uncertain{maximal} \ill{} \ill{} \ill{}
\uncertain{engendré} \add{\ill{}} \ill{}. \ill{} \ill{} \ill{}

\struck{\ill{}} \emph{Cor.}\ \struck{\ill{}} \ill{} \ill{}
$B = A[t_{1}, \ldots, t_{n}]$ \ill{} \ill{} $= \mathfrak{m}$, \uncertain{idéal}
\uncertain{maximal}
\[
\dim B_{\mathfrak{m}} = \dim A + n
\]
\emph{Cor.} \uncertain{Si} $A$ \struck{\ill{}} \ill{},
$\dim A[t_{1}, \ldots, t_{n}] = \dim A + n$

\emph{Cor.} \ill{} \ill{} \ill{} \ill{} \ill{} [\ill{} \ill{} \ill{} \ill{}
\uncertain{de} \uncertain{type} \uncertain{fini} \ill{} $k$ \ill{} \ill{}
\uncertain{caractéristique}] \ill{} \ill{}, \ill{} \ill{} \ill{}
\[
\dim B_{\mathfrak{m}} = \dim A + n - \deg \mathrm{tr}\, k(\mathfrak{m})/k(\mathfrak{m}')
\]
\[
\nu = \dim \bigl[B_{\mathfrak{m}} \otimes k = k[t_{1}, \ldots, t_{n}]_{\mathfrak{m}}\bigr]
\]
\note{le dernier indice $\mathfrak{m}'$ est une lecture ; $\nu$ est accroché
par une accolade à « $n - \deg \mathrm{tr}$ ».}

\struck{\emph{Cor.} $A$ \ill{} \ill{}. \uncertain{Alors} $\dim A[t_{1},
\ldots, t_{n}] = \dim A + n$.}

\struck{(ii) \uncertain{Si} $A = k[t_{1}, \ldots, t_{n}]_{\mathfrak{m}}$,
\ill{} \ill{}, \ill{} \ill{} \ill{} \ill{} $A$ \ill{} \ill{} \ill{}}
\note{deux lignes barrées horizontalement, reliées par un cadre à une
reprise en dessous.}
\[
\mathrm{prof}\, B_{\mathfrak{m}} = \mathrm{prof}\, A + \nu \qquad
\mathrm{coprof}\, B_{\mathfrak{m}} = \mathrm{coprof}\, A
\]
\note{devant la première égalité, un « $A$ » et un mot biffé
(\uncertain{« Cor. »}) ; les lectures « prof », « coprof » sont incertaines.}

\emph{Cor.} $X = Y[t_{1}, \ldots, t_{n}]$, $Y$ \uncertain{loc.}\
\uncertain{noeth.} \uncertain{Soit} $x \in X$, $y \in Y$. \uncertain{Soit}
$P$ \ill{} \ill{} \uncertain{propriétés} \ldots\ \uncertain{Alors} $X$
\ill{} $P$ \ill{} \ill{} \uncertain{sur} $Y$ \ill{} \ill{} \ill{} $y$.
\ill{} \ill{} \ill{} \uncertain{propriétés}).

\page{129}
\note{un grand crochet et deux traits obliques barrent le haut de la page,
jusqu'à la ligne de l'$\mathrm{Ext}^{-2}$ ; lecture incertaine de la prose.}
\uncertain{Soit} \ill{} \uncertain{élément} $\xi \in \mathrm{Ext}^{i}(C; G, S\gamma)$,
i.e.\ un \uncertain{homom.}\ \struck{\uncertain{de} \uncertain{degré}}
\ill{}
\[
G \to S\gamma,
\]
d'\uncertain{où} un \uncertain{homom.}\ \ill{}
\[
Rf_{*}G \to Rf_{*}S\gamma,
\]
\uncertain{donc} \uncertain{il} \uncertain{reste} à \uncertain{définir} un
\uncertain{homom.}\ \ill{} \uncertain{degré} \ill{}
\[
Rf_{*}S\gamma = Rf_{*}\bigl(\gamma \otimes \mu_{n}(1)\bigr) \longrightarrow \gamma
\]
\note{sous « $\otimes \mu_{n}$ », un petit « $R^{\cdot}$ » avec une marque
en exposant, peu lisible.}
i.e.\ un \uncertain{élément} \struck{de}
\struck{$\mathrm{Ext}^{-2}\bigl(k, Rf_{*}(\gamma \otimes \mu_{n}), \gamma\bigr)$} \ldots
\ill{}, \ill{} \uncertain{homom.}\ \uncertain{évident}
\[
\underline{E}_{f}(G, S\gamma) \longrightarrow \underline{E}_{f}(Rf_{*}G, Rf_{*}S\gamma)
\]
et on \uncertain{vient} \uncertain{de} \uncertain{définir} \ill{}
\uncertain{élément}
\[
\gamma \in H^{0}(Rf_{*}S\gamma, \gamma)
\]
\note{le symbole devant la parenthèse est biffé et remplacé au-dessus
par un $H$ \uncertain{d'exposant $0$}, souligné par une vague.}
\uncertain{d'où} \ill{} \uncertain{définissons} un \struck{\ill{}}
\uncertain{homom.}
\[
\underline{E}_{f}(Rf_{*}G, Rf_{*}S\gamma) \longrightarrow \underline{E}_{f}(Rf_{*}G, \gamma)
\]
\ill{} \ill{} \uncertain{composant}, on \uncertain{trouve}
\uncertain{l'homom.}\ \uncertain{voulu}
\[
\underline{E}_{f}(G, S\gamma) \longrightarrow \underline{E}_{f}(Rf_{*}G, \gamma)
\]
\uncertain{Ceci} \uncertain{dit}, le \uncertain{th.}\ \uncertain{de}
\uncertain{dualité} \uncertain{s'énonce} \add{\ill{}
\struck{\uncertain{faisceaux} \uncertain{algébriques}}
\uncertain{faisceaux} \uncertain{de} Weil}

\emph{Th.} (i) \uncertain{Si} $G$ \ill{} \uncertain{fini}
[\uncertain{ou} \uncertain{plus} \uncertain{généralement} \ill{} \ill{}
$H^{i}(G) \neq 0$ \ill{} \struck{\ill{}}], \struck{\uncertain{alors}}
\uncertain{plus} \uncertain{généralement} \ill{} \uncertain{affine}, et
\ill{} \uncertain{composante} \uncertain{abélienne}, \uncertain{alors}
\[
\underline{E}_{f}(G, S\gamma) \xrightarrow{\ \sim\ } \underline{E}(Rf_{*}G, \gamma)
\]
(ii) \uncertain{Si} $G$ \uncertain{est} \add{\uncertain{de} \uncertain{type}
\uncertain{fini}} \uncertain{sans} \uncertain{composante}
\uncertain{abélienne}, \uncertain{alors}
\[
\underline{E}_{f}(G, S\gamma) \longrightarrow \underline{E}(Rf_{*}G, \gamma) \ \ldots
\]
\ill{} \uncertain{mod} $\Theta$ [\uncertain{i.e.} les $\underline{E}$
\uncertain{dans} $\mathbb{Q}/\mathbb{Z}$ \struck{\ill{}} \ill{} \ill{} \ill{}
\ill{} \ldots].
\marginal{\uncertain{précisions}}
\note{« de type fini », en (ii), est inscrit dans une ellipse au-dessus de
« sans composante » ; un trait vertical borde la fin du (i) ; le mot de
marge, écrit en biais à gauche du (ii), n'est relié à rien.}

\page{130}
\note{feuillet vide, retourné ; en haut, à la plume, une en-tête annulée
par trois traits obliques.}
$\underline{\mathrm{V}}$ \struck{\ill{}}

\struck{\emph{Schémas de Hilbert, de Picard}}

\struck{(\uncertain{sans} \uncertain{théorie} des \uncertain{schémas}
\uncertain{abéliens} \ill{} \ill{})}
\note{« Schémas de Hilbert, de Picard » est souligné deux fois.}

\page{131}
On \uncertain{a} une \uncertain{suite} \uncertain{exacte} de
$C$-\uncertain{groupes}
\[
0 \to \mu_{n} \to \mathbb{G}_{m} \xrightarrow{\ n\ } \mathbb{G}_{m} \to 0
\]
d'où
\[
H^{1}(C, \mathbb{G}_{m}) \xrightarrow{\ n\ } H^{1}(C, \mathbb{G}_{m}) \to H^{2}(C, \mu_{n}) \to H^{2}(C, \mathbb{G}_{m})
\]
[N.B. On \uncertain{a} $\underline{H}^{0}_{f}(\mathbb{G}_{m}) \simeq \mathbb{G}_{m}$,
$\underline{H}^{i}_{f}(\mathbb{G}_{m}) = 0$ pour $i \geqslant 2$, d'où une
suite exacte \ldots]
\[
0 \to H^{1}(k, \mathbb{G}_{m}) \to H^{1}(C, \mathbb{G}_{m}) \to H^{0}\bigl(k, \underline{H}^{1}_{f}(\mathbb{G}_{m})\bigr) \to H^{2}(k, \mathbb{G}_{m})
\]
\note{sous $H^{1}(k, \mathbb{G}_{m})$ : « $= 0$ », avec la glose
\uncertain{« th. 90 »} ; sous $H^{0}(k, \underline{H}^{1}_{f}(\mathbb{G}_{m}))$ :
« $= \mathrm{Pic}(C/k)$ » ; sous $H^{2}(k, \mathbb{G}_{m})$ :
« $= \mathrm{Br}(k)$ ».}

\uncertain{Ainsi}, $H^{1}(C, \mathbb{G}_{m}) \subset \mathrm{Pic}\, C$

$\gamma$ \uncertain{groupe} \add{\uncertain{fini}}
\struck{\uncertain{algébrique}} sur $k$, \uncertain{définissons} un
\uncertain{groupe} \struck{$\gamma \otimes \mu_{n}$ \ill{} $k$, \ill{}}
$= D(\Delta\gamma)$ de type \uncertain{multiplicatif}
[\uncertain{isomorphe} à $\gamma \otimes \mu_{n}$ si $n\gamma = 0$].
\uncertain{Posons} $S\gamma = (D\Delta\gamma)(1)$
\note{la torsion finale est lue $(1)$, comme en batch 5 (page 99) ; le
chiffre pourrait être un $2$.}
\uncertain{On} \uncertain{va} \uncertain{définir} un \struck{\ill{}}
\uncertain{accouplement}
\begin{itemize}
\item[a)] $\mathrm{Ext}^{i}(C; G, S\gamma) \longrightarrow \mathrm{Ext}^{i\ldots}(k, Rf_{*}G, \gamma)$
\end{itemize}
\note{l'exposant du second $\mathrm{Ext}$ est surchargé : un $i$ suivi d'un
trait épais illisible (peut-être $i-2$).}
\uncertain{associé} à un \uncertain{homom.}
\begin{itemize}
\item[b)] (*) $\underline{E}_{f}(G, S\gamma) \longrightarrow \underline{E}_{k}(Rf_{*}G, \gamma)$
\end{itemize}
\note{l'astérisque est entouré.}
[\uncertain{évidemment} \uncertain{fonctoriel} \uncertain{en} $G$, $\gamma$
\ldots], \uncertain{on} \uncertain{définit} \uncertain{aussi} un
\uncertain{homom.}
\[
\underline{\mathrm{Ext}}^{i}_{f}(G, S\gamma) \longrightarrow \underline{\mathrm{Ext}}^{i}_{k}(Rf_{*}G, \gamma)
\]
\begin{itemize}
\item[c)]
\end{itemize}
\note{« c) » reste sans suite.}

\page{132}
\note{la page est un tableau à deux colonnes, séparées par un trait vertical,
et quatre étages séparés par des traits horizontaux, avec une troisième
colonne étroite à droite pour les deux premiers étages ; on le donne étage
par étage.}
\begin{itemize}
\item Premier étage. À gauche : $\mu_{n}$, $(n,p) = 1$ ; puis
  $\mu_{n}$, ${}_{n}J$, $\mathbb{Z}/n\mathbb{Z}$. À droite :
  $\mathbb{Z}/n\mathbb{Z}$, $(n,p) = 1$ ; puis
  $\mathbb{Z}/n\mathbb{Z} = \mu_{n} \otimes \check{\mu}_{n}$,
  ${}_{n}J \otimes \check{\mu}_{n}$, $\check{\mu}_{n}$. Dans la colonne
  étroite : « \uncertain{autodualité} \struck{\ill{}} \uncertain{ens.}\
  \uncertain{des} \uncertain{pts} \uncertain{d'ordre} $n$ de la
  \uncertain{jacobienne} », suivi d'un mot barré d'un trait épais.
\item Deuxième étage. À gauche : $\mu_{p}$ ; puis $\mu_{p}$, ${}_{p}J$,
  $\mathbb{Z}/p\mathbb{Z}$. À droite : $\mathbb{Z}/p\mathbb{Z}$ ; puis
  $\mathbb{Z}/p\mathbb{Z}$, ${}_{(1-F)}H^{1}(C, \underline{\mathcal{O}}_{C})^{F}$,
  $H^{1}(C, \underline{\mathcal{O}}_{C})_{(1-F)} = 0$. Dans la colonne
  étroite :
  \[
  \Delta({}_{p}J) \simeq H^{1}(C, \underline{\mathcal{O}}_{C})^{F} \quad \bigl(H^{1}(C, \mathbb{Z}/p\mathbb{Z})\bigr)
  \]
  « i.e.\ » $\pi_{1}(C) \otimes \mathbb{Z}/p\mathbb{Z} \simeq {}_{p}J$
  [\uncertain{autodualité} \uncertain{de} la \uncertain{jacobienne}
  \ill{} $V_{p}$ ?].
\item Troisième étage. À gauche : $\alpha_{p}$ ; puis $\alpha_{p}$,
  ${}_{F}H^{1}(C, \underline{\mathcal{O}}_{C})$,
  $H^{1}(C, \underline{\mathcal{O}}_{C})_{F}$ ; et [\uncertain{provient}
  \uncertain{des} \uncertain{complexes} \emph{\uncertain{relatifs}} \ill{}
  $A$ \uncertain{par} \ill{} \add{\uncertain{Frobenius}}
  \struck{\ill{}}]. À droite : $\alpha_{p}$ ; puis $0$ \uncertain{(?)}. Dans
  un cadre à droite : « Dualité \struck{\ill{}} \ill{} de
  $H^{1}(C, \underline{\mathcal{O}}_{C})_{F}$ et
  ${}_{F}H^{1}(C, \underline{\mathcal{O}}_{C})$ provient de
  l'\uncertain{autodualité} \uncertain{d'une} \uncertain{jacobienne},
  \ill{} \ill{} \ill{} $A$ \ill{} \ill{} \ill{} \ill{} ${}_{p}A$) \ill{} :
  \ill{} \ill{}, \ill{} \ill{} \ill{} \ill{} \ill{} \ill{} ${}_{p}A$ \ill{}
  \ill{} \uncertain{propriétés} \ill{} \uncertain{dualité} \uncertain{dualité}
  \uncertain{provenant} \uncertain{d'une} \uncertain{dualité} \ill{}
  \uncertain{correspondantes} \ill{} \uncertain{algèbres} \ill{} ».
\item Quatrième étage. À gauche : $\mathbb{G}_{m}$ ; puis
  $\mathbb{G}_{m}$, $\underline{\mathrm{Pic}}^{\cdot}_{C/k}$, $0$. À droite :
  $\mathbb{Z}$ ; puis
  \[
  \mathbb{Z},\ 0,\ H^{2}(C, \mathbb{Q}/\mathbb{Z}) = \mathrm{Hom}\bigl(T.(J), \mathbb{Q}/\mathbb{Z}\bigr),\ H^{3}(C, \mathbb{Q}/\mathbb{Z}) \simeq \varinjlim_{(n,p)=1} \check{\mu}_{n},\ 0
  \]
\end{itemize}
\marginal{\uncertain{provient} \uncertain{des} \uncertain{complexes}
\uncertain{identiques} ; $J$ \uncertain{par} \ill{} \ill{}}
\note{note écrite verticalement dans la marge gauche, en face des deux
premiers étages, reliée par une accolade à la colonne de gauche.}

\uncertain{Pour} $G = \mathbb{G}_{m}$, le \uncertain{th.}\ \uncertain{de}
\uncertain{dualité} \uncertain{s'écrit} \struck{\ill{}}
\[
H^{*}(C, \gamma) \simeq \underline{\underline{\mathrm{Ext}}}^{*}\bigl(Rf_{*}(\mathbb{G}_{m})(-1), \gamma\bigr) \Longleftarrow \underline{\underline{\mathrm{Ext}}}^{p+q}\bigl(R^{-q}f_{*}(\mathbb{G}_{m}), \gamma\bigr)
\]
\note{sous $Rf_{*}(\mathbb{G}_{m})$, une accolade et la glose
« \uncertain{complexe} \uncertain{idéal} \ill{} ] ». L'exposant de droite
est lu $p+q$.}
et on \uncertain{se} \uncertain{ramène} \uncertain{au} \uncertain{cas}
\emph{\uncertain{géométrique}}, \ill{} $\gamma = \mathbb{Q}/\mathbb{Z}$ :
\[
H^{*}(C, \mathbb{Q}/\mathbb{Z}) \simeq \underline{\underline{\mathrm{Ext}}}^{*}\bigl(Rf_{*}(\mathbb{G}_{m})(1), \gamma\bigr) \Longleftarrow \underline{\underline{\mathrm{Ext}}}^{p-q}\bigl(R^{-q}f_{*}(\mathbb{G}_{m}), \mathbb{Q}/\mathbb{Z}\bigr).
\]
\note{la torsion est ici $(1)$, là-haut $(-1)$ ; l'exposant de droite est
lu $p-q$ sans certitude.}
et \uncertain{on} \uncertain{a} \uncertain{évidemment} :
\[
H^{0}(C, \mathbb{Q}/\mathbb{Z}) \simeq \mathbb{Q}/\mathbb{Z}, \quad H^{1}(C, \mathbb{Q}/\mathbb{Z}) \simeq \mathrm{Ext}^{1}(J, \mathbb{Q}/\mathbb{Z}), \quad H^{2}(C, \mathbb{Q}/\mathbb{Z}) \simeq \ill{} = \check{\mu}
\]
\note{au-dessus de $\mathrm{Ext}^{1}(J, \mathbb{Q}/\mathbb{Z})$, reliée par
« $=$ » : « $\mathrm{Ext}^{1}(J^{0}, \mathbb{Q}/\mathbb{Z})$ » ; au-dessous,
sous une flèche : « \uncertain{autodualité} \uncertain{de}
\uncertain{la} \uncertain{jacobienne} ».}

\page{133}
\note{feuillet presque vide ; en haut, une en-tête annulée par cinq traits
obliques.}
$\underline{\mathrm{V}}$

\struck{\emph{Fidèle platitude}}

\struck{\emph{Descente} fid.\ plate}

\struck{Applications}

\page{134}
\uncertain{Ci-dessus}, on a \uncertain{négligé} un \uncertain{coefficient}
\ill{} \ill{} \add{\uncertain{affirmations}} \emph{\uncertain{provenant}
\uncertain{de} la} \ill{}

\uncertain{Regardons} le \uncertain{cas} \uncertain{où} $F$
\uncertain{faisceau} de Weil \emph{\uncertain{concentré}} en un
\uncertain{point} $x \in C$. \uncertain{Alors} \struck{\ill{}}
[\uncertain{on} \ill{} \ill{}]
\[
T(G) \simeq G_{x}
\]
« considéré comme groupe algébrique sur $k$ »
\note{le $G$ de gauche est surchargé ; au-dessus de « considéré », un petit
signe d'insertion sans mot.}
\uncertain{On} \uncertain{doit} \uncertain{prouver}
\[
\mathrm{Ext}^{i}(C; G, \gamma \otimes \mu_{n}) \simeq \mathrm{Ext}^{i-2}(k; G_{x}, \gamma)
\]
\uncertain{Je} \uncertain{ne} \uncertain{suis} \uncertain{pas}
\uncertain{encore} \uncertain{sûr} \uncertain{que} le $\mathrm{Ext}^{i}$
\uncertain{écrit} \uncertain{est} \uncertain{aussi} le $\mathrm{Ext}^{i}$
\struck{\uncertain{local}} \uncertain{formé} \ill{} \ill{} \uncertain{complété}
de $\mathcal{O}_{C,x}$ \ldots\ [\uncertain{si} $G = F_{Y}$ (\uncertain{où}
$Y = \lbrace x \rbrace$)],
\marginal{$Y$ \uncertain{partie} \uncertain{finie} \uncertain{de} $C$}
\note{note de marge entourée, en face des deux lignes précédentes.}
\uncertain{Pour} le \uncertain{cas} \uncertain{général} \uncertain{des}
$\mathrm{Ext}$, \ill{} \ill{} \ill{} \ill{} \ill{} \ill{} \uncertain{où}
$i \colon U \to C$ \ill{} \uncertain{l'immersion} \ill{}
$(U = C - Y)$ \ill{}
\struck{$\mathrm{Ext}^{*}(C; G, \gamma \otimes \mu_{n}) \Longleftarrow H^{p}\bigl(C, \underline{\mathrm{Ext}}^{q}(G, \gamma \otimes \mu_{n})\bigr)$}
\ill{} \ill{} \ill{} \ill{}, \ill{} \ill{} \ill{} \ill{} $C(F)$, [i.e.\ :
\uncertain{prouver} \ldots]
\[
\mathrm{Ext}^{i}(U; F, S\gamma) \simeq \mathrm{Ext}^{i}_{\ill{}}\bigl(k; Rf_{U*}(i_{!}F), \gamma\bigr)
\]
\note{au-dessus de $i_{!}F$ : « $F_{U}$ », relié par un arc ; l'indice de
$f$ est surchargé.}
\uncertain{où} $f_{U} \colon U \to \mathrm{Spec}(k)$, \ill{} \ill{}
\uncertain{prouver} \ldots
\struck{($F$ \uncertain{faisceau} \uncertain{plat} \ill{} $U$)}
\struck{$F$} \ill{}. \ill{} \ill{} \ill{} à \uncertain{prouver} \uncertain{que}
\[
\boxed{H^{i}(U; \hat{F}) \simeq \mathrm{Ext}^{i-2}\bigl(Rf_{*}(i_{!}F_{U}), \mathbb{Q}/\mathbb{Z}\bigr)}
\]
\note{l'exposant de $H$ est surchargé ; l'argument de $Rf_{*}$ aussi,
lu $i_{!}F_{U}$.}
\struck{\uncertain{N.B.} \ill{} \uncertain{distinguer}} \struck{\ill{}}
\ill{} $G$ \uncertain{On} \uncertain{a} \uncertain{défini} \ill{}
\uncertain{suite} \uncertain{exacte} (\ill{} \ill{}
\struck{$\ill{}$}\ $F_{Y} \to \ldots$) :
\[
\ldots \longleftarrow H^{i}(U; \hat{F}) \longleftarrow \mathrm{Ext}^{i-2}\bigl(Rf_{*}(F_{!}), \mathbb{Q}/\mathbb{Z}\bigr) \longleftarrow \mathrm{Ext}^{i-2}(Rf_{*}F_{Y}, \mathbb{Q}/\mathbb{Z}) \longleftarrow
\]
\note{au-dessus du premier terme, un « $\mathrm{Ext}$ » biffé.}
\uncertain{On} \uncertain{le} \uncertain{vérifie} \ill{} \uncertain{un}
\uncertain{calcul} \uncertain{facile} \struck{\uncertain{Correctible}}
(\uncertain{en} \uncertain{utilisant} \ill{} \ill{} \uncertain{résultats}
\uncertain{de} \uncertain{dualité} \ill{} \uncertain{vérifiés} \ldots)
\uncertain{dans} le \uncertain{cas} \ill{} $F = \mu_{n}$, \uncertain{d'où}
$\hat{F} = \mathbb{Z}/n\mathbb{Z}$ [\uncertain{d'un} \ill{}], (\ill{}
\add{\uncertain{premier} \uncertain{à} $p$ \ldots} \struck{\ill{}}), \ill{}
\uncertain{récurrence}, \struck{\ill{}}.

\page{135}
\note{au dos d'une page dactylographiée qui transparaît.}
\[
F = i_{*}(G)
\]
\[
\left\lbrace
\begin{aligned}
&H^{0}_{x}(F) = H^{1}_{x}(F) = 0 \\
&H^{2}_{x}(F) = H^{1}(K, F)
\end{aligned}
\right.
\]
\[
\underline{\underline{\mathrm{Ext}}}^{i}\bigl(H^{1}(K, F), \mathbb{Q}/\mathbb{Z}\bigr) \simeq \mathrm{Ext}^{i}\bigl(S, i_{*}(G), \ill{} \otimes \mathbb{G}_{m}\bigr)
\]
\note{sous le membre de droite, un « $\mathrm{Ext}^{i}$ » isolé. Le
$\mathbb{Q}/\mathbb{Z}$ de gauche est surchargé.}
\[
H^{i}_{x}(F) \to H^{i}(X, F) \to H^{i}(U, F)
\]
\[
\longleftarrow H^{3}(U, \hat{F})
\]
\note{au-dessus de $H^{3}(U, \hat{F})$, un $\mathrm{Ext}$ biffé.}
\[
\underline{H}_{x}(F) \to \underline{H}_{X}(F) \to \underline{H}_{U}(F)
\]
\note{un arc revient de $\underline{H}_{U}(F)$ vers $\underline{H}_{x}(F)$ ;
au-dessus de la ligne, un « $F$ ». À droite : « \uncertain{triangle}
\uncertain{exact} \uncertain{de} \uncertain{complexes} \uncertain{de}
$k$-\uncertain{groupes} ».}
\uncertain{Prenons} \uncertain{les} \uncertain{duaux} \uncertain{dans}
$\mathbb{Q}/\mathbb{Z}$
\[
E(X; F, \mu_{\infty}) \longleftarrow E_{x}(X; F, \mu_{\infty}) \longleftarrow E(U; F, \mu_{\infty})
\]
\note{un arc revient du premier terme vers le dernier.}
\[
\left\lbrace
\begin{aligned}
E(U; F, \mu_{\infty})^{n} &\simeq E\bigl(k, \underline{H}_{U}(F), \mathbb{Q}/\mathbb{Z}\bigr)^{n-1} \\
E(X; F, \mu_{\infty})^{n} &\simeq E\bigl(k, \underline{H}_{X}(F), \mathbb{Q}/\mathbb{Z}\bigr)^{n-2} \\
E_{x}(X; F, \mu_{\infty})^{n} &\simeq E\bigl(k, \underline{H}_{x}(F), \mathbb{Q}/\mathbb{Z}\bigr)^{n-2}
\end{aligned}
\right.
\]
\note{dans un grand cadre, avec un grand « ? » à gauche. Entre les lignes,
des flèches verticales marquées $\partial$ et $i$, à gauche entre les $E$
et à droite entre les $\underline{H}$.}
\[
\begin{array}{ll}
E_{U}(X; F, \mu_{\infty})^{i} & E\bigl(k, \underline{H}_{U}(F), \mathbb{Q}/\mathbb{Z}\bigr)^{i-1} \\
E_{X}(X; F, \mu_{\infty}) & E\bigl(k, \underline{H}_{X}(F), \mathbb{Q}/\mathbb{Z}\bigr) \\
E_{x}(X; F, \mu_{\infty})^{i} & E\bigl(k, \underline{H}_{x}(F), \mathbb{Q}/\mathbb{Z}\bigr)^{i-2}
\end{array}
\]
\note{deuxième tableau dans le même cadre, sans signe entre les colonnes.
Colonne de droite de la page, séparée par un trait vertical :}
\[
E(U; F, \mu_{\infty})^{n} \simeq E(X, F_{U}, \mu_{\infty})^{n}
\overset{\wr}{\simeq} E\bigl(k, \underline{H}_{X}(F_{U}), \mathbb{Q}/\mathbb{Z}\bigr)^{n-2}
\overset{?}{\simeq} E\bigl(k, \underline{H}_{U}(F), \mathbb{Q}/\mathbb{Z}\bigr)^{n-1}
\]
\[
\left\lbrace
\begin{aligned}
&\underline{H}^{m}_{X}(F_{U})^{n} \simeq \underline{H}^{m}_{x}(F)^{n} \\
&\boxed{\underline{H}^{m}_{U}(F) \simeq \underline{H}^{m+1}_{x}(F_{U})}
\end{aligned}
\right.
\]
\note{la seconde ligne est encadrée, avec un « ? » dans le cadre ; les
exposants $n$ de la première ligne sont peu sûrs.}
\[
E_{x}(X; F, \mu_{\infty})^{n} \simeq E(X, F_{x}, \mu_{\infty})^{n}
\overset{\wr}{\simeq} E\bigl(k, \underline{H}_{X}(F_{x}), \mathbb{Q}/\mathbb{Z}\bigr)^{n-2}
\overset{?}{\simeq} E\bigl(k, \underline{H}_{x}(F), \mathbb{Q}/\mathbb{Z}\bigr)^{n-2}
\]
\note{dans la colonne de droite, les isomorphismes sont écrits
verticalement, « $\wr$ » puis « $\wr$ ? », les termes empilés.}

\page{136}
\note{feuillet dactylographié paginé « - 37 - », « n° 338 », non
recomposé : la fin d'un énoncé sur un anneau $B$ et un homomorphisme
$h \colon A \to B$ (conditions (iv) : $B$ plat à droite et
$h^{-1}(B\mathfrak{a}) = \mathfrak{a}$ pour tout idéal à gauche
$\mathfrak{a}$ ; (v) : $B$ plat et tout idéal à gauche maximal
$\mathfrak{m}$ de $A$ est de la forme $h^{-1}(\mathfrak{m}')$), la
démonstration (i) $\Rightarrow$ (iii) $\Rightarrow$ (iv) $\Rightarrow$ (ii)
$\Rightarrow$ (i) de la caractérisation de la fidèle platitude, puis
l'énoncé de la Proposition 9 commencé. Seule intervention à la main, au
crayon dans la marge gauche, en biais, en face de (v) :
« $\mathrm{Spec}(B) \to \mathrm{Spec}(A)$ \uncertain{surjectif} », avec un
trait vertical.}

\page{137}
\note{au dos d'une page dactylographiée qui transparaît. Dans un grand
cadre, deux colonnes de termes reliés par des flèches verticales, la
gauche descendante, la droite montante sauf sa première flèche, qui descend ; on les redessine en grille.}

\begin{tikzcd}[column sep=small, row sep=small, nodes={font=\scriptsize}]
\underline{H}^{i-1}_{S/k}(G) \arrow[d, no head] & \underline{H}^{2-(i-1)}_{Y/k}(S, \hat{G}) \arrow[d] \\
\underline{H}^{i-1}_{U/k}(G) \arrow[d] & \underline{H}^{1-(i-1)}_{U/k}(\hat{G}) \\
\underline{H}^{i}_{Y/k}(S, G) \arrow[d] & \underline{H}^{2-i}_{S/k}(\hat{G}) \arrow[u, Rightarrow] \\
\underline{H}^{i}_{S/k}(G) \arrow[d] & \underline{H}^{2-i}_{Y/k}(\hat{G}) \arrow[u] \\
\underline{H}^{i}_{U/k}(G) & \underline{H}^{1-i}_{U/k}(\hat{G}) \arrow[u]
\end{tikzcd}

\note{la première flèche de gauche est un simple trait. Au-dessus de
l'exposant $2-(i-1)$ est écrit « $3-i$ », au-dessus de $1-(i-1)$
« $2-i$ » ; aux deux termes du milieu de la colonne de droite, un
« $\mathrm{Ext}$ » est biffé devant le $H$. La flèche montante entre
eux est grasse. En haut à droite, hors du cadre :
« $\underline{H}^{i}_{S/k}(F)$ » et un « $H$ » isolé.}

Sous le cadre :
\struck{$\underline{H}^{i}_{U/k}(G)$} \qquad \struck{\ill{}}
\[
\overset{?}{\simeq}\ \underline{H}^{i+1}_{S/k}(\tilde{G})
\]
\note{le second terme est surchargé ; son exposant est lu $i+1$ sans
certitude. Plus bas, un « $H$ » isolé.}

\page{139}
\note{au dos d'une page dactylographiée qui transparaît.}
\struck{$H^{i}(C,$}
\[
\underline{H}^{i}_{U/k}(F) \qquad \mathrm{Ext}^{i}
\]
\[
\underline{\mathrm{Ext}}^{i}_{U/k}(G, \tilde{\gamma}) \simeq \underline{\mathrm{Ext}}^{i}_{k}\bigl(\mathcal{N}^{\cdot}(G), \gamma\bigr)
\]
\note{devant $(G)$, dans l'argument de droite, une lettre biffée.}
\[
\begin{array}{cc}
\underline{\mathrm{Ext}}^{i}_{S/k}(G, \tilde{\gamma}) & \underline{\mathrm{Ext}}^{i}_{k}\bigl(\mathcal{N}^{\cdot}(G), \gamma\bigr) \\
\downarrow & \uparrow \\
\underline{\mathrm{Ext}}^{i}_{U/k}(G, \tilde{\gamma}) & \underline{\mathrm{Ext}}^{i}_{k}\bigl(\mathcal{N}^{\cdot}(G_{U}), \gamma\bigr)
\end{array}
\]
\note{l'indice du premier terme est surchargé ($U$ récrit en $S$ ?).}
\[
0 \to G_{U} \to G \to G^{Y} \to 0
\]
\[
\underline{\mathrm{Ext}}^{*}_{U/k}(G, \tilde{\gamma}) = \underline{\mathrm{Ext}}^{*}_{S/k}(G_{U}, \tilde{\gamma}) = \underline{\mathrm{Ext}}^{*}_{k}\bigl(\mathcal{N}^{\cdot}(i_{!}(G)), \gamma\bigr)
\]
\note{sous $G_{U}$, « $= i_{!}(G_{U})$ » écrit verticalement ; sous le
dernier terme, une flèche vers la gauche.}
\[
\underline{H}^{*}_{U/k}(\hat{G})
\]
\note{écrit sous le premier terme, relié par « $=$ » vertical.}
\[
\boxed{\underline{H}^{*}_{Y/k}\bigl(i_{!}(G)\bigr) \overset{?}{\simeq} \underline{H}^{*}_{U/k}(G)} \quad ???
\]
\[
\underline{H}^{*}_{Y/k}\bigl(i_{!}(G)\bigr) \simeq \ill{}
\]
\note{le dernier signe ressemble à un $\vdash$.}

\end{document}
